% !TeX root = ../../main.tex
% sections/chap3/sec_model_risk_sharing.tex

\section{リスク共有}
\label{sec:sec_model_risk_sharing}

前節では家計\(h\)の最適化問題を定式化した。
本節ではモデルの仮定である国内完備市場が家計間のリスク共有にどのような含意をもたらすかを説明する。
これにより、なぜ所得の限界効用、消費、そして可処分所得が家計間で共通化されるのかを数学的に明らかにする。

\subsection{自国家計のリスク共有}
\label{sec:sec_model_risk_sharing_home}

\subsubsection{所得の限界効用\(\lambda_t^h\)の共通化}
\label{sec:sec_model_risk_sharing_home_lambda_h}
本モデルでは自国と外国のそれぞれにおいて国内の債券市場は完備であると仮定しており、
家計はあらゆる不確実性に対応した国内コンティンジェント債券\(d_{t+1}^{h\to H}\)を取引できる。
任意の2つの家計\(h\)と\(h'\)を考える。
両者ともすべての時点・すべての状態で最適化行動をとるため
国内コンティンジェント債券の購入に関するFOC（FOC）は常に成立する。
\begin{equation}
\label{form:sec_model_risk_sharing_home_lambda_h_eq}
    \lambda_t^hq_{t,t+1}=\beta_t^H\pi\lambda_{t+1}^h
\end{equation}
\begin{equation}
\label{form:sec_model_risk_sharing_home_lambda_h2_eq}
    \lambda_t^{h'}q_{t,t+1}=\beta_t^H\pi\lambda_{t+1}^{h'}
\end{equation}
ここで\(q_{t,t+1}\)は状態コンティンジェント債の価格であり、\(\pi\)は状態の発生確率である。
これらの比をとることで家計間の限界効用の比率が時間や状態に依存しない定数\(k\)となることがわかる。
\begin{equation}
\label{form:sec_model_risk_sharing_home_lambda_h_ratio_eq}
    \frac{\lambda_t^h}{\lambda_t^{h'}}=\frac{\lambda_{t+1}^h}{\lambda_{t+1}^{h'}}=k
\end{equation}
この定数\(k\)の値を特定するため、ショックの発生時点\(s\)に注目する。
すべての家計が同一の効用関数をもち、かつ初期状態において資産保有量がゼロ（\(d_s^h=b_s^h=0\)）であることから、
各家計が直面する最適化問題は数学的に完全に同一となる。
したがって、その解として得られる所得の限界効用もすべての家計で一致しなければならない
（\(\lambda_s^h=\lambda_s^{h'}\)）。
この条件を式\eqref{form:sec_model_risk_sharing_home_lambda_h_ratio_eq}に適用することで、\(k=1\)が特定される。
一度\(k=1\)が確定すれば、国内債券市場が完備であるかぎり、
将来にわたっていかなるショックが発生してもすべての家計の所得の限界効用は常に一致し続ける。
\begin{equation}
\label{form:sec_model_risk_sharing_home_lambda_h_equality_eq}
    \lambda_t^h=\lambda_t^{h'}\quad（\text{for all }h,h'\in H,\text{ all }t）
\end{equation}

\subsubsection{総消費指数\(c_t^{h\to W}\)の共通化}
\label{sec:sec_model_risk_sharing_home_c_h_W}
消費に関するFOC\eqref{form:model_households_indices_aggregation_stage2_home_lambda_h_p_H_W_c_h_W_eq}を
総消費指数\(c_t^{h\to W}\)について解けば\(c_t^{h\to W}=1/（p_t^{H\to W}\lambda_t^h）\)となる。
これと式\eqref{form:sec_model_risk_sharing_lambda_h_equality}を用いると以下の通り
すべての家計の総消費指数が一致することがわかる。
\begin{equation}
\label{form:sec_model_risk_sharing_home_c_h_W_equality_eq}
    c_t^{h\to W}
    =\frac{1}{p_t^{H\to W}\lambda_t^h}
    =\frac{1}{p_t^{H\to W}\lambda_t^{h'}}
    =c_t^{h'\to W}
    \quad（\text{for all }h,h'\in H,\text{ all }t）
\end{equation}
こうしてすべての家計の総消費指数は同一の値をとることが示された。

\subsubsection{可処分所得の共通化}
\label{sec:sec_model_risk_sharing_home_omega_h_equality}
まず、家計\(h\)が第\(t\)期において実際に利用可能な名目資源の総計を名目可処分所得\(\omega_t^h\)と定義する。
\begin{equation}
\label{form:sec_model_risk_sharing_omega_def}
\omega_t^h\equiv d_t^h+(1+i_{t-1}^F)\frac{e_t^{/*}}{e_{t-1}^{/*}}b_t^h+(1-\tau_t^H)p_t^hy_t^h+t_t^H
\end{equation}
ここで、ある時点\(t=s\)において、家計間で名目可処分所得に差（\(\omega_s^{h'}>\omega_s^{h''}\)）があると仮定する。
すべての家計の総消費指数\(c_s^{h\to W}\)は共通であるため、
家計\(h'\)の名目可処分所得の超過分\(\omega_s^{h'}-\omega_s^{h''}>0\)はすべて蓄積に充てられなければならない。
さらに将来のすべての期\(t>s\)においても総消費指数\(c_t^{h\to W}\)は全家計で共通である。
このため、家計\(h'\)が第\(s\)期に蓄積した名目的な余剰資源が将来の消費支出に振り向けられることはなく、
この資産価値の差は利息を伴って将来にわたって持ち越され続けることとなる。
しかし家計の効用関数は消費に関して単調増加であり、
家計は与えられた予算制約の下で生涯効用を最大化するように行動している。
将来にわたって一切消費に充てられない名目資産を保有し続けることはその資源をいずれかの期の消費に回すことで
生涯効用をさらに高められる余地があることを意味しており、これは家計の効用最大化行動と矛盾する。
ゆえに初期の仮定は誤りであり、すべての時点・状態において名目可処分所得は家計間で一致しなければならない。
\begin{equation}
\label{form:sec_model_risk_sharing_omega_h_equality}
\omega_t^h=\omega_t^{h'}\quad（\text{for all }h,h'\in H,\text{ all }t）
\end{equation}

\subsubsection{自国財消費指数\(c_t^{h\to H}\)の共通化}
\label{sec:sec_model_risk_sharing_home_c_h_H}
自国財消費指数\(c_t^{h\to H}\)への需要関数は以下のようであった
\begin{equation}
\tag{\ref{form:model_households_indices_aggregation_stage1b_home_c_h_H_eq}}
    c_t^{h\to H}=\alpha^H\frac{p_t^{H\to W}}{p_t^H}c_t^{h\to W}
\end{equation}
この右辺を構成するすべての変数
（選好パラメータ\(\alpha^H\)、物価指数\(p_t^{H\to W}, p_t^H\)、総消費指数\(c_t^{h\to W}\)）
は家計間で共通であるため、左辺の\(c_t^{h\to H}\)もまたすべての家計間で共通となる。
\begin{equation}
\label{form:sec_model_risk_sharing_home_c_h_H_equality_eq}
    c_t^{h\to H}=c_t^{h'\to H}\quad（\text{for all }h,h'\in H,\text{ all }t）
\end{equation}

\subsubsection{外国財消費指数\(c_t^{h\to F}\)の共通化}
\label{sec:sec_model_risk_sharing_home_c_h_F}
同様に外国財消費指数\(c_t^{h\to F}\)への需要関数は以下の通りであった。
\begin{equation}
\tag{\ref{form:model_households_indices_aggregation_stage1b_home_c_h_F_eq}}
    c_t^{h\to F}=（1-\alpha^H）\frac{p_t^{H\to W}}{p_t^F}c_t^{h\to W}
\end{equation}
この右辺を構成するすべての変数
（選好パラメータ\(\alpha^H\)、物価指数\(p_t^{H\to W}, p_t^F\)、総消費指数\(c_t^{h\to W}\)）
が家計間で共通であるため、左辺の\(c_t^{h\to F}\)もまたすべての家計間で共通となる。
\begin{equation}
\label{form:sec_model_risk_sharing_home_c_h_F_equality_eq}
    c_t^{h\to F}=c_t^{h'\to F}\quad（\text{for all }h,h'\in H,\text{ all }t）
\end{equation}

\subsection{外国家計のリスク共有}
\label{sec:sec_model_risk_sharing_foreign}
外国家計についても自国家計と同様に、国内債券市場が完備であるとの仮定から家計間での完全なリスク共有が実現される。
まず外国家計\(f\)の外国通貨建て名目可処分所得\(\omega_t^{f*}\)を以下の通り定義する。
\begin{equation}
\label{form:sec_model_risk_sharing_foreign_omega_star_def}
\omega_t^{f*}\equiv d_t^{f\to F*}+（1+i_{t-1}^H）\frac{e_{t-1}^{f*}}{e_t^{f*}}b_t^{f\to H*}+（1-\tau_t^F）p_t^{f*}y_t^f+t_t^{F*}
\end{equation}
このとき自国家計と同様の最適化条件および背理法による議論によりすべての時点・状態において、
外国家計間における所得の限界効用\(\lambda_t^{f/*}\)、総消費指数\(c_t^{f\to W}\)、
および名目可処分所得\(\omega_t^{f*}\)は以下の通り共通化される。
\begin{equation}
\label{form:sec_model_risk_sharing_foreign_lambda_f_slash_star_equality}
\lambda_t^{f/*}=\lambda_t^{f'/*}\quad（\text{for all }f,f'\in F,\text{ all }t）
\end{equation}
\begin{equation}
\label{form:sec_model_risk_sharing_foreign_c_f_W_equality}
c_t^{f\to W}=c_t^{f'\to W}\quad（\text{for all }f,f'\in F,\text{ all }t）
\end{equation}
\begin{equation}
\label{form:sec_model_risk_sharing_foreign_omega_star_equality}
\omega_t^{f*}=\omega_t^{f'/*}=\omega_t^{F/*}\quad（\text{for all }f,f'\in F,\text{ all }t）
\end{equation}
さらにこれらに伴い外国財消費指数\(c_t^{f\to F}\)および自国財消費指数\(c_t^{f\to H}\)についても
すべての外国家計で共通の値をとる。
\begin{equation}
\label{form:sec_model_risk_sharing_foreign_c_f_F_equality}
c_t^{f\to F}=c_t^{f'\to F}\quad（\text{for all }f,f'\in F,\text{ all }t）
\end{equation}
\begin{equation}
\label{form:sec_model_risk_sharing_foreign_c_f_H_equality}
c_t^{f\to H}=c_t^{f'\to H}\quad（\text{for all }f,f'\in F,\text{ all }t）
\end{equation}